% ===========================================================================
% Código Sintético — cabecera LaTeX del INTERIOR IMPRESO
% Formato: 6 x 9 pulgadas (15,24 x 22,86 cm), márgenes espejados
% Uso: pandoc --include-in-header=tools/latex/header-interior.tex
% La marca @RUNNING-TITLE@ la sustituye tools/build_editions.sh por el
% título de la serie/volumen que aparece en las páginas pares (verso).
% ===========================================================================

% --- Imágenes: solo se reducen, nunca desbordan la caja de texto -----------
% (adjustbox aporta la clave "max width", que actúa únicamente si la imagen
%  excede la medida; evita ampliar diagramas pequeños.)
\usepackage[export]{adjustbox}
\let\interiorincludegraphics\includegraphics
\renewcommand{\includegraphics}[2][]{%
  \adjustbox{max width=\linewidth,max totalheight=0.88\textheight}{%
    \interiorincludegraphics[#1]{#2}}}

% --- Raíz de los recursos gráficos del manuscrito --------------------------
% Los capítulos referencian las imágenes como "imagenes/..."; se declara la
% raíz para que LaTeX las encuentre desde el directorio de compilación.
\graphicspath{{@GRAPHICSPATH@}}

% --- Símbolos: tipografía de respaldo para signos fuera de Pagella ---------
% Comprobado con una compilación de prueba: TeX Gyre Pagella sí tiene → ■ ▼ ✓
% – — … € ≈ ≥ ≤ ≠ ⊂ ⌊ ⌋ − y griego; pero NO tiene subíndices, ► ni ✗.
% Sin este respaldo, esos signos saldrían en blanco en el papel.
% Evita el conflicto entre las tablas de protrusión y caracteres activos.
\microtypesetup{protrusion=false}
\usepackage{newunicodechar}
\usepackage{fontspec}
\newfontfamily\symfont{DejaVu Sans}
\newcommand{\sym}[1]{{\symfont #1}}
% Subíndices (CO₂, R², aₙ, xᵢ, S₋₁…)
\newunicodechar{₀}{\sym{₀}}
\newunicodechar{₁}{\sym{₁}}
\newunicodechar{₂}{\sym{₂}}
\newunicodechar{₃}{\sym{₃}}
\newunicodechar{₄}{\sym{₄}}
\newunicodechar{₅}{\sym{₅}}
\newunicodechar{₆}{\sym{₆}}
\newunicodechar{₇}{\sym{₇}}
\newunicodechar{₈}{\sym{₈}}
\newunicodechar{₉}{\sym{₉}}
\newunicodechar{₋}{\sym{₋}}
\newunicodechar{ₙ}{\sym{ₙ}}
\newunicodechar{ᵢ}{\sym{ᵢ}}
% Señalética de los capítulos
\newunicodechar{►}{\sym{►}}
\newunicodechar{✗}{\sym{✗}}
% Equivalentes de impresión; el EPUB conserva los caracteres originales.
\newunicodechar{～}{\textasciitilde}
\newunicodechar{👍}{aprobación}
\newunicodechar{👎}{desaprobación}

% --- Código: cuerpo menor y líneas partidas --------------------------------
% Pandoc define el entorno `Highlighting` sobre fancyvrb SIN ajuste de línea;
% en una caja de 4,5" eso desbordaría la mancha y la guillotina cortaría el
% código. Se recarga con fvextra y se activa el corte (incluso a mitad de
% palabra), que es lo correcto para bloques de código en formato de libro.
\usepackage{fvextra}
\fvset{breaklines=true, breakanywhere=true, fontsize=\footnotesize,
       breaksymbolleft={}, breaksymbolright={}}
\DefineVerbatimEnvironment{Highlighting}{Verbatim}{breaklines,breakanywhere,commandchars=\\\{\}}
% Pandoc usa verbatim para bloques sin lenguaje; también deben partir líneas.
\DefineVerbatimEnvironment{verbatim}{Verbatim}{breaklines,breakanywhere,fontsize=\footnotesize}
% Identificadores largos en \texttt: punto de corte tras cada guion bajo.
% (pandoc escapa "_" como "\_".) seqsplit se descarta: rompe con comandos
% dentro de \texttt (\^{}, \textbackslash...), comprobado en vol3 cap. 23.
\let\interiortexttt\texttt
\renewcommand{\texttt}[1]{\begingroup
  \def\_{\char`\_\allowbreak}%
  \interiortexttt{#1}\endgroup}

% --- Tablas: cuerpo menor y separación silábica en columnas estrechas -----
\usepackage{ragged2e}
\let\raggedright\RaggedRight
\setlength{\LTleft}{0pt}
\setlength{\LTright}{0pt}
\setlength{\LTcapwidth}{\linewidth}
\renewcommand{\arraystretch}{1.15}
\AtBeginEnvironment{longtable}{\small\setlength{\tabcolsep}{3pt}}

% --- Cabeceras y pies profesionales: folio al margen exterior -------------
\usepackage{fancyhdr}
\renewcommand{\chaptermark}[1]{\markboth{#1}{}}
\pagestyle{fancy}
\fancyhf{}
\fancyhead[LE]{\small\itshape @RUNNING-TITLE@}
\fancyhead[RO]{\small\itshape\leftmark}
\fancyfoot[LE,RO]{\small\thepage}
\renewcommand{\headrulewidth}{0.4pt}
\renewcommand{\footrulewidth}{0pt}
% Páginas de apertura de capítulo: folio abajo, sin filete ni cabecera
\fancypagestyle{plain}{%
  \fancyhf{}%
  \fancyfoot[LE,RO]{\small\thepage}%
  \renewcommand{\headrulewidth}{0pt}%
}
\fancypagestyle{empty}{%
  \fancyhf{}%
  \renewcommand{\headrulewidth}{0pt}%
}

% --- Control de viudas y huérfanas -----------------------------------------
% (microtype ya lo carga la plantilla de pandoc: no se repite aquí.)
\widowpenalty=1000
\clubpenalty=1000
\displaywidowpenalty=1000
\raggedbottom